\documentclass[11pt]{article}
\newcommand{\ListaTitulo}{Lista 05 --- Submuestreo e Efeitos Aleatórios}
% Preâmbulo compartilhado — MATD48, Listas de Exercícios 2026
% Usado por Lista{NN}.tex e Gabarito{NN}.tex via % Preâmbulo compartilhado — MATD48, Listas de Exercícios 2026
% Usado por Lista{NN}.tex e Gabarito{NN}.tex via % Preâmbulo compartilhado — MATD48, Listas de Exercícios 2026
% Usado por Lista{NN}.tex e Gabarito{NN}.tex via \input{preamble}

\usepackage[utf8]{inputenc}
\usepackage[T1]{fontenc}
\usepackage[brazil]{babel}
\usepackage{fullpage}
\usepackage{amsmath,amssymb}
\usepackage{enumitem}
\usepackage{xcolor}
\usepackage{fancyhdr}
\usepackage{titlesec}
\usepackage{listings}
\usepackage{hyperref}
\usepackage{booktabs}
\usepackage{graphicx}

\definecolor{matd48azul}{HTML}{0B4F6C}
\definecolor{matd48verde}{HTML}{2E8B57}

\titleformat{\section}{\normalfont\Large\bfseries\color{matd48azul}}{\thesection}{1em}{}
\titleformat{\subsection}{\normalfont\large\bfseries\color{matd48azul}}{\thesubsection}{1em}{}

\providecommand{\ListaTitulo}{Lista de Exercícios}

\setlength{\headheight}{14pt}
\pagestyle{fancy}
\fancyhf{}
\lhead{\small MATD48 --- Planejamento de Experimentos A}
\rhead{\small \ListaTitulo}
\cfoot{\thepage}
\renewcommand{\headrulewidth}{0.4pt}

\lstset{
  basicstyle=\ttfamily\small,
  breaklines=true,
  frame=single,
  columns=fullflexible,
  backgroundcolor=\color{gray!8},
  commentstyle=\color{matd48verde},
  keywordstyle=\color{matd48azul}\bfseries,
  language=R,
  extendedchars=true,
  literate=%
    {á}{{\'a}}1 {à}{{\`a}}1 {â}{{\^a}}1 {ã}{{\~a}}1
    {é}{{\'e}}1 {ê}{{\^e}}1 {í}{{\'i}}1 {ó}{{\'o}}1
    {ô}{{\^o}}1 {õ}{{\~o}}1 {ú}{{\'u}}1 {ç}{{\c c}}1
    {Á}{{\'A}}1 {À}{{\`A}}1 {Â}{{\^A}}1 {Ã}{{\~A}}1
    {É}{{\'E}}1 {Ê}{{\^E}}1 {Í}{{\'I}}1 {Ó}{{\'O}}1
    {Ô}{{\^O}}1 {Õ}{{\~O}}1 {Ú}{{\'U}}1 {Ç}{{\c C}}1
}

\hypersetup{colorlinks=true, linkcolor=matd48azul, urlcolor=matd48azul, citecolor=matd48azul}

\newcommand{\cabecalho}[2]{%
  \begin{center}
    {\Large\bfseries\color{matd48azul} #1}\\[0.3em]
    {\large #2}\\[0.3em]
    \rule{\linewidth}{0.6pt}
  \end{center}
  \vspace{0.5em}
}

\newenvironment{questao}[1]{\vspace{1em}\noindent\textbf{Questão #1.}\hspace{0.4em}}{\par}
\newenvironment{solucao}{\vspace{0.3em}\noindent\textit{Solução.}\hspace{0.4em}\color{black!85}}{\par\vspace{0.5em}}


\usepackage[utf8]{inputenc}
\usepackage[T1]{fontenc}
\usepackage[brazil]{babel}
\usepackage{fullpage}
\usepackage{amsmath,amssymb}
\usepackage{enumitem}
\usepackage{xcolor}
\usepackage{fancyhdr}
\usepackage{titlesec}
\usepackage{listings}
\usepackage{hyperref}
\usepackage{booktabs}
\usepackage{graphicx}

\definecolor{matd48azul}{HTML}{0B4F6C}
\definecolor{matd48verde}{HTML}{2E8B57}

\titleformat{\section}{\normalfont\Large\bfseries\color{matd48azul}}{\thesection}{1em}{}
\titleformat{\subsection}{\normalfont\large\bfseries\color{matd48azul}}{\thesubsection}{1em}{}

\providecommand{\ListaTitulo}{Lista de Exercícios}

\setlength{\headheight}{14pt}
\pagestyle{fancy}
\fancyhf{}
\lhead{\small MATD48 --- Planejamento de Experimentos A}
\rhead{\small \ListaTitulo}
\cfoot{\thepage}
\renewcommand{\headrulewidth}{0.4pt}

\lstset{
  basicstyle=\ttfamily\small,
  breaklines=true,
  frame=single,
  columns=fullflexible,
  backgroundcolor=\color{gray!8},
  commentstyle=\color{matd48verde},
  keywordstyle=\color{matd48azul}\bfseries,
  language=R,
  extendedchars=true,
  literate=%
    {á}{{\'a}}1 {à}{{\`a}}1 {â}{{\^a}}1 {ã}{{\~a}}1
    {é}{{\'e}}1 {ê}{{\^e}}1 {í}{{\'i}}1 {ó}{{\'o}}1
    {ô}{{\^o}}1 {õ}{{\~o}}1 {ú}{{\'u}}1 {ç}{{\c c}}1
    {Á}{{\'A}}1 {À}{{\`A}}1 {Â}{{\^A}}1 {Ã}{{\~A}}1
    {É}{{\'E}}1 {Ê}{{\^E}}1 {Í}{{\'I}}1 {Ó}{{\'O}}1
    {Ô}{{\^O}}1 {Õ}{{\~O}}1 {Ú}{{\'U}}1 {Ç}{{\c C}}1
}

\hypersetup{colorlinks=true, linkcolor=matd48azul, urlcolor=matd48azul, citecolor=matd48azul}

\newcommand{\cabecalho}[2]{%
  \begin{center}
    {\Large\bfseries\color{matd48azul} #1}\\[0.3em]
    {\large #2}\\[0.3em]
    \rule{\linewidth}{0.6pt}
  \end{center}
  \vspace{0.5em}
}

\newenvironment{questao}[1]{\vspace{1em}\noindent\textbf{Questão #1.}\hspace{0.4em}}{\par}
\newenvironment{solucao}{\vspace{0.3em}\noindent\textit{Solução.}\hspace{0.4em}\color{black!85}}{\par\vspace{0.5em}}


\usepackage[utf8]{inputenc}
\usepackage[T1]{fontenc}
\usepackage[brazil]{babel}
\usepackage{fullpage}
\usepackage{amsmath,amssymb}
\usepackage{enumitem}
\usepackage{xcolor}
\usepackage{fancyhdr}
\usepackage{titlesec}
\usepackage{listings}
\usepackage{hyperref}
\usepackage{booktabs}
\usepackage{graphicx}

\definecolor{matd48azul}{HTML}{0B4F6C}
\definecolor{matd48verde}{HTML}{2E8B57}

\titleformat{\section}{\normalfont\Large\bfseries\color{matd48azul}}{\thesection}{1em}{}
\titleformat{\subsection}{\normalfont\large\bfseries\color{matd48azul}}{\thesubsection}{1em}{}

\providecommand{\ListaTitulo}{Lista de Exercícios}

\setlength{\headheight}{14pt}
\pagestyle{fancy}
\fancyhf{}
\lhead{\small MATD48 --- Planejamento de Experimentos A}
\rhead{\small \ListaTitulo}
\cfoot{\thepage}
\renewcommand{\headrulewidth}{0.4pt}

\lstset{
  basicstyle=\ttfamily\small,
  breaklines=true,
  frame=single,
  columns=fullflexible,
  backgroundcolor=\color{gray!8},
  commentstyle=\color{matd48verde},
  keywordstyle=\color{matd48azul}\bfseries,
  language=R,
  extendedchars=true,
  literate=%
    {á}{{\'a}}1 {à}{{\`a}}1 {â}{{\^a}}1 {ã}{{\~a}}1
    {é}{{\'e}}1 {ê}{{\^e}}1 {í}{{\'i}}1 {ó}{{\'o}}1
    {ô}{{\^o}}1 {õ}{{\~o}}1 {ú}{{\'u}}1 {ç}{{\c c}}1
    {Á}{{\'A}}1 {À}{{\`A}}1 {Â}{{\^A}}1 {Ã}{{\~A}}1
    {É}{{\'E}}1 {Ê}{{\^E}}1 {Í}{{\'I}}1 {Ó}{{\'O}}1
    {Ô}{{\^O}}1 {Õ}{{\~O}}1 {Ú}{{\'U}}1 {Ç}{{\c C}}1
}

\hypersetup{colorlinks=true, linkcolor=matd48azul, urlcolor=matd48azul, citecolor=matd48azul}

\newcommand{\cabecalho}[2]{%
  \begin{center}
    {\Large\bfseries\color{matd48azul} #1}\\[0.3em]
    {\large #2}\\[0.3em]
    \rule{\linewidth}{0.6pt}
  \end{center}
  \vspace{0.5em}
}

\newenvironment{questao}[1]{\vspace{1em}\noindent\textbf{Questão #1.}\hspace{0.4em}}{\par}
\newenvironment{solucao}{\vspace{0.3em}\noindent\textit{Solução.}\hspace{0.4em}\color{black!85}}{\par\vspace{0.5em}}


\begin{document}

\cabecalho{Lista de Exercícios 05}{DCA com submuestreo e modelo de efeitos aleatórios (Capítulo 3 / Aula 05)}

\noindent \textit{Entrega: até a próxima aula. Justifique todas as respostas --- nestas listas, a
justificativa vale mais que a resposta final. O gabarito completo está em \texttt{Gabarito05.pdf}.}

\begin{questao}{1} \textbf{(Agricultura/pecuária --- identificação de unidades com submuestreo)}

Um zootecnista testa o efeito de três dietas (A, B, C) sobre o teor de gordura no leite (\%).
Doze vacas são usadas, sorteadas 4 para cada dieta; cada vaca é ordenhada 3 vezes ao longo do
dia, e o teor de gordura é medido em cada ordenha separadamente.

\begin{enumerate}[label=(\alph*)]
  \item Identifique a unidade experimental e a unidade amostral. Quantos níveis $t$, UEs por
    tratamento $r$ e submuestras por UE $q$ há neste desenho? Confira que $N = trq$ bate com o
    número total de medições.
  \item Escreva o modelo estatístico apropriado para $Y_{ijk}$, definindo cada termo.
  \item Se o zootecnista rodasse uma ANOVA de um fator diretamente nas $N$ medições de gordura
    (ignorando que vêm de apenas 12 vacas), qual seria o erro cometido e por que ele tende a
    produzir significância "fácil demais"?
\end{enumerate}
\end{questao}

\begin{questao}{2} \textbf{(Dedução --- esperança do quadrado médio do erro experimental)}

Considere o modelo balanceado de submuestreo $Y_{ijk} = \mu + \tau_i + \delta_{ij} +
\varepsilon_{ijk}$, $i=1,\ldots,t$, $j=1,\ldots,r$, $k=1,\ldots,q$, com $\delta_{ij}
\overset{iid}{\sim} N(0,\sigma_\delta^2)$ e $\varepsilon_{ijk} \overset{iid}{\sim}
N(0,\sigma_\varepsilon^2)$ independentes entre si. A soma de quadrados do erro experimental é
$$
SQ_{\text{exp}} = q\sum_{i=1}^t\sum_{j=1}^r (\bar Y_{ij\cdot} - \bar Y_{i\cdot\cdot})^2 .
$$
Você pode usar, sem demonstrar, o fato de que se $X_1,\ldots,X_r$ são i.i.d. com variância
$\sigma_X^2$, então $E\!\left[\sum_{j=1}^r (X_j-\bar X)^2\right] = (r-1)\sigma_X^2$.

\begin{enumerate}[label=(\alph*)]
  \item Mostre que $\bar Y_{ij\cdot} - \bar Y_{i\cdot\cdot} = (\delta_{ij}-\bar\delta_{i\cdot}) +
    (\bar\varepsilon_{ij\cdot} - \bar\varepsilon_{i\cdot\cdot})$, e dê a variância de
    $\delta_{ij}+\bar\varepsilon_{ij\cdot}$ (trate as $r$ quantidades $\delta_{ij} +
    \bar\varepsilon_{ij\cdot}$, $j=1,\ldots,r$, como i.i.d. dentro de cada tratamento $i$).
  \item Usando o item (a) e o fato fornecido, mostre que $E[SQ_{\text{exp}}] = t(r-1)(\sigma_\varepsilon^2
    + q\sigma_\delta^2)$.
  \item Conclua que $E[MS_{\text{exp}}] = E[SQ_{\text{exp}}]/[t(r-1)] = \sigma_\varepsilon^2 +
    q\sigma_\delta^2$, confirmando a coluna $E[QM]$ da tabela de ANOVA com submuestreo.
\end{enumerate}
\end{questao}

\begin{questao}{3} \textbf{(Agricultura/pecuária --- ANOVA com submuestreo, tabela incompleta)}

Retomando o experimento das dietas (Questão 1: $t=3$, $r=4$, $q=3$, $N=36$), o ajuste da ANOVA
aninhada produziu os valores abaixo (alguns propositalmente omitidos, marcados com "?").

\begin{center}
\begin{tabular}{lccccc}
\toprule
Fonte & g.l. & SQ & QM & $F$ \\
\midrule
Dieta & ? & 54{,}0 & ? & ? \\
Erro experimental (entre vacas) & ? & 39{,}0 & ? & --- \\
Erro de amostragem (entre ordenhas) & ? & 48{,}0 & 2{,}00 & --- \\
Total & 35 & ? & & \\
\bottomrule
\end{tabular}
\end{center}

\begin{enumerate}[label=(\alph*)]
  \item Complete todos os graus de liberdade, o $SQ_{\text{Total}}$, e os quadrados médios
    faltantes.
  \item Calcule o $F$ \textbf{correto} para testar o efeito da dieta. Usando
    $F_{2,9;0{,}05}\approx 4{,}26$, há evidência de efeito de dieta?
  \item Calcule o que aconteceria se, por engano, o zootecnista usasse o erro de amostragem no
    denominador em vez do erro experimental. Compare esse $F$ "errado" (e seus graus de liberdade)
    com o $F$ correto do item (b), e explique em uma frase por que o erro é sempre nessa direção
    (favorece rejeitar $H_0$ com mais facilidade do que deveria).
\end{enumerate}
\end{questao}

\begin{questao}{4} \textbf{(Psicologia --- componentes de variância)}

Um pesquisador sorteia 6 vídeos motivacionais de um banco de 40 vídeos possíveis e mostra cada
um a 10 novos participantes antes de um teste de resistência física, medindo o tempo de
resistência (segundos). O objetivo não é comparar esses 6 vídeos específicos, mas estimar o
quanto a escolha do vídeo, de modo geral, contribui para a variabilidade do desempenho. A ANOVA
fornece $MS_{\text{Entre vídeos}} = 850$ e $MS_E = 310$.

\begin{enumerate}[label=(\alph*)]
  \item Qual modelo (efeitos fixos ou aleatórios) é apropriado aqui? Justifique com base no
    processo de seleção dos vídeos.
  \item Estime $\hat\sigma_\tau^2$ (variância entre vídeos) e $\hat\sigma^2$ (variância residual).
  \item Calcule o coeficiente de correlação intraclasse (ICC) e interprete seu valor em uma frase.
\end{enumerate}
\end{questao}

\begin{questao}{5} \textbf{(Dedução --- variância e correlação no modelo de efeitos aleatórios)}

No modelo de efeitos aleatórios $Y_{ij} = \mu + \tau_i + \varepsilon_{ij}$, $i=1,\ldots,t$,
$j=1,\ldots,n$, com $\tau_i \overset{iid}{\sim} N(0,\sigma_\tau^2)$ independente de
$\varepsilon_{ij}\overset{iid}{\sim} N(0,\sigma^2)$:

\begin{enumerate}[label=(\alph*)]
  \item Obtenha $\mathrm{Var}(Y_{ij})$.
  \item Obtenha $\mathrm{Cov}(Y_{ij}, Y_{ij'})$ para $j\neq j'$ (mesmo nível $i$) e
    $\mathrm{Cov}(Y_{ij}, Y_{i'j'})$ para $i \neq i'$ (níveis diferentes).
  \item Usando os itens (a) e (b), mostre que $\mathrm{Corr}(Y_{ij},Y_{ij'}) =
    \sigma_\tau^2/(\sigma_\tau^2+\sigma^2)$ para $j\neq j'$, e explique por que essa quantidade é
    chamada de coeficiente de correlação intraclasse.
\end{enumerate}
\end{questao}

\end{document}
